\documentclass[12pt]{article}
\usepackage{amsmath, amssymb}
\usepackage[margin=1in]{geometry}
\setlength{\parskip}{6pt}
\title{QUBO Formulation: Graph K-Coloring Problem}
\date{}
\begin{document}
\maketitle

\subsection*{Graph K-Coloring Problem}

\paragraph{Problem Statement.}
Given an undirected graph $G=(V,E)$ with $N=|V|$ vertices and a set of $K$ available colors, assign exactly one color to each vertex such that no two adjacent vertices share the same color. The goal is to find a valid $K$-coloring if one exists, or otherwise minimize the number of monochromatic (violated) edges. A valid coloring corresponds to an objective value of zero.

\paragraph{Decision Variables.}
For each vertex $v \in V$ and each color $c \in \{0, 1, \dots, K-1\}$, introduce a binary decision variable $x_{v,c} \in \{0,1\}$. The variable $x_{v,c}$ equals $1$ if and only if vertex $v$ is assigned color $c$, and $0$ otherwise.

\paragraph{QUBO Derivation.}
We construct a Quadratic Unconstrained Binary Optimization (QUBO) formulation by relaxing the hard constraints into the objective function via quadratic penalty terms.

\textbf{Step 1 -- Objective Function.}
The original goal is to minimize the number of violated edges. An edge $(u,v) \in E$ is violated if both endpoints receive the same color $c$. This is captured by the product $x_{u,c}x_{v,c}$. The objective to be minimized is:
\begin{align}
    \min \quad \sum_{(u,v) \in E} \sum_{c=0}^{K-1} x_{u,c} x_{v,c}.
\end{align}

\textbf{Step 2 -- Constraints.}
Each vertex must be assigned exactly one color. For every vertex $v \in V$, this hard constraint is expressed as:
\begin{align}
    \sum_{c=0}^{K-1} x_{v,c} = 1, \quad \forall v \in V.
\end{align}

\textbf{Step 3 -- Penalty Term Expansion.}
We enforce the constraints by adding a quadratic penalty to the objective. Let $\lambda > 0$ be the penalty weight. The penalty contribution for a single vertex $v$ is:
\begin{align}
    P_v(\mathbf{x}) = \lambda \left( \sum_{c=0}^{K-1} x_{v,c} - 1 \right)^2.
\end{align}
Expanding the square yields:
\begin{align}
    \left( \sum_{c=0}^{K-1} x_{v,c} - 1 \right)^2 = \sum_{c=0}^{K-1} x_{v,c}^2 + \sum_{c=0}^{K-1} \sum_{\substack{c'=0 \\ c' \neq c}}^{K-1} x_{v,c} x_{v,c'} - 2 \sum_{c=0}^{K-1} x_{v,c} + 1.
\end{align}
Using the idempotent property of binary variables, $x_{v,c}^2 = x_{v,c}$, we simplify:
\begin{align}
    \left( \sum_{c=0}^{K-1} x_{v,c} - 1 \right)^2 = \sum_{c=0}^{K-1} x_{v,c} + \sum_{c \neq c'} x_{v,c} x_{v,c'} - 2 \sum_{c=0}^{K-1} x_{v,c} + 1 = \sum_{c \neq c'} x_{v,c} x_{v,c'} - \sum_{c=0}^{K-1} x_{v,c} + 1.
\end{align}
Summing over all vertices, the total penalty contribution to the Hamiltonian is:
\begin{align}
    P(\mathbf{x}) = \lambda \sum_{v \in V} \left( \sum_{c \neq c'} x_{v,c} x_{v,c'} - \sum_{c=0}^{K-1} x_{v,c} + 1 \right).
\end{align}

\textbf{Step 4 -- Combined Hamiltonian.}
Adding the objective and penalty terms gives the full QUBO Hamiltonian $H(\mathbf{x})$:
\begin{align}
    H(\mathbf{x}) = \sum_{(u,v) \in E} \sum_{c=0}^{K-1} x_{u,c} x_{v,c} + \lambda \sum_{v \in V} \left( \sum_{c \neq c'} x_{v,c} x_{v,c'} - \sum_{c=0}^{K-1} x_{v,c} + 1 \right).
\end{align}

\textbf{Step 5 -- Q Matrix Entries and Offset.}
We express $H(\mathbf{x})$ in the standard QUBO form $H(\mathbf{x}) = \sum_{i} Q_{ii} x_i + \sum_{i < j} 2 Q_{ij} x_i x_j + c$, where $i$ indexes the pair $(v,c)$. Reading off the coefficients from the expanded Hamiltonian yields the following closed-form expressions for the Q matrix entries and constant offset:
\begin{align}
    Q_{(v,c),(v,c)} &= -\lambda, \quad \forall v \in V, \, c \in \{0,\dots,K-1\}, \\
    Q_{(v,c),(v,c')} &= \lambda, \quad \forall v \in V, \, c \neq c', \\
    Q_{(u,c),(v,c)} &= \frac{1}{2}, \quad \forall (u,v) \in E, \, c \in \{0,\dots,K-1\}, \\
    c &= N \lambda.
\end{align}
All other entries $Q_{ij}$ are zero.

\paragraph{Penalty Weight Justification.}
Let $\Delta = \max_{v \in V} \deg(v)$ denote the maximum vertex degree in $G$. Each variable $x_{v,c}$ participates in at most $\deg(v)$ terms in the objective function. Consequently, flipping a single variable can reduce the objective value by at most $\Delta$. To guarantee that any constraint violation is strictly more costly than any potential objective improvement, we require $\lambda > \Delta$. Setting $\lambda$ to any value strictly greater than $\Delta$ (e.g., $\lambda = \Delta + 1$) ensures that infeasible solutions are energetically unfavorable compared to all feasible solutions.

\paragraph{Correctness Argument.}
Minimizing $H(\mathbf{x})$ over $\{0,1\}^{NK}$ recovers the optimal coloring because (i) any feasible assignment satisfies $\sum_{c} x_{v,c} = 1$ for all $v$, causing all penalty terms to vanish and reducing $H(\mathbf{x})$ exactly to the objective function counting violated edges; (ii) any infeasible assignment incurs a penalty of at least $\lambda$ for each violated constraint. Since $\lambda > \Delta$, the penalty strictly exceeds the maximum possible decrease in the objective function achievable by any single variable flip. Therefore, the global minimum of $H(\mathbf{x})$ must occur at a feasible point, and among feasible points, $H(\mathbf{x})$ is minimized precisely when the number of violated edges is minimized.

\paragraph{Illustrative Example.}
Consider a concrete instance with $N=3$ vertices, $K=3$ colors, and a complete graph $G=K_3$ where every pair of vertices is connected. The edges are $(0,1)$, $(1,2)$, and $(2,0)$. The maximum degree is $\Delta=2$, so we set $\lambda=10$. The variable mapping is $i = vK + c$, yielding indices $0,1,2$ for vertex $0$, $3,4,5$ for vertex $1$, and $6,7,8$ for vertex $2$.

The concrete objective counts monochromatic edges:
\begin{align}
    \text{Obj}(\mathbf{x}) = \sum_{c=0}^{2} \left( x_{0,c}x_{1,c} + x_{1,c}x_{2,c} + x_{2,c}x_{0,c} \right).
\end{align}
The concrete penalty terms for each vertex $v \in \{0,1,2\}$ are:
\begin{align}
    P_v(\mathbf{x}) = 10 \left( \sum_{c \neq c'} x_{v,c}x_{v,c'} - \sum_{c=0}^{2} x_{v,c} + 1 \right).
\end{align}
The resulting Hamiltonian $H(\mathbf{x})$ is the sum of the objective and all three penalty terms. The non-zero Q matrix entries for this instance are:
\begin{align}
    Q_{0,0} = Q_{1,1} = Q_{2,2} = Q_{3,3} = Q_{4,4} = Q_{5,5} = Q_{6,6} = Q_{7,7} = Q_{8,8} &= -10, \\
    Q_{0,1} = Q_{0,2} = Q_{1,2} = Q_{3,4} = Q_{3,5} = Q_{4,5} = Q_{6,7} = Q_{6,8} = Q_{7,8} &= 10, \\
    Q_{0,3} = Q_{1,4} = Q_{2,5} = Q_{3,6} = Q_{4,7} = Q_{5,8} = Q_{6,0} = Q_{7,1} = Q_{8,2} &= \frac{1}{2},
\end{align}
with offset $c = 3 \times 10 = 30$.

A valid 3-coloring exists. Setting $x_{0,0}=1$, $x_{1,1}=1$, $x_{2,2}=1$ (all other variables $0$) satisfies the exactly-one-color constraints and yields zero violated edges. The optimal energy is $H(\mathbf{x}^*) = 0 + 30 = 30$. Any assignment violating the constraints incurs a penalty of at least $10$, strictly exceeding the maximum possible objective reduction of $2$, confirming the formulation's correctness for this instance.

\end{document}